% Generated by render.py; edit preamble.tex, citations.json and routes.tikz
\documentclass[border=12pt]{standalone}
\def\pgfsysdriver{pgfsys-dvisvgm.def}
\usepackage[dvisvgm]{xcolor}
\usepackage{xeCJK}
\setCJKmainfont{HaranoAjiGothic-Regular.otf}[BoldFont=HaranoAjiGothic-Bold.otf]
\setCJKsansfont{HaranoAjiGothic-Regular.otf}[BoldFont=HaranoAjiGothic-Bold.otf]
\usepackage{amsmath,amssymb,tikz,lmodern}
\usetikzlibrary{arrows.meta,positioning,calc,fit,backgrounds}
\renewcommand{\familydefault}{\sfdefault}
\definecolor{ink}{HTML}{25322F}
\definecolor{green}{HTML}{246657}
\definecolor{tint}{HTML}{EEF4F0}
\definecolor{soft}{HTML}{C6D2CC}
\definecolor{muted}{HTML}{52615B}
\newif\ifJapanese
\newcommand{\Tx}[2]{\ifJapanese #1\else #2\fi}
\newcommand{\Eff}{\overline{\operatorname{Eff}}}
\newcommand{\heading}[3]{{\color{green}\bfseries #1\quad\Tx{#2}{#3}}\\[5pt]}
% A white rectangle makes the complete citation label clickable, including spaces.
\newcommand{\refurl}[2]{\special{dvisvgm:raw <a xlink:href="#1" target="_blank" rel="noopener noreferrer">}{\setlength{\fboxsep}{0pt}\colorbox{white}{\color{green}\underline{\strut #2}}}\special{dvisvgm:raw </a>}}
\tikzset{
 card/.style={draw=soft,fill=white,rounded corners=2pt,align=left,text=ink,font=\fontsize{10.5}{14}\selectfont,inner xsep=12pt,inner ysep=10pt,text width=14.4cm},
 half/.style={card,text width=7.5cm,minimum height=2.05cm},
 result/.style={card,draw=green,fill=tint,line width=.9pt},
 reason/.style={align=left,text=muted,font=\fontsize{9}{12}\selectfont,inner sep=0pt,text width=12.6cm},
 branchreason/.style={reason,text width=6.9cm},
 edge/.style={-{Stealth[length=2.4mm]},draw=green,line width=1pt},
 wire/.style={draw=green,line width=1pt},
}
% Keep the arrow in a dedicated left gutter; explanations never cover it.
\def\flow#1#2#3;{
 \coordinate (#2-from) at ($(#1.south)+(-6.9,0)$);
 \coordinate (#2-to) at ($(#2.north)+(-6.9,0)$);
 \draw[edge] (#2-from)--(#2-to);
 \node[reason,anchor=west] (#2-reason) at ($(#2-from)!.5!(#2-to)+(.55,0)$) {#3};
}

% Style adapted from the repository's Schnell diagrams; MIT attribution retained.

% Generated from citations.json; do not edit.
\newcommand{\Cite}[1]{\csname cite-#1\endcsname}
\expandafter\def\csname cite-catalog\endcsname{\shortstack[l]{\refurl{https://github.com/openai/math/blob/1581d4cbdd69a0faf34b2089614fe2a814316a6a/overview.pdf}{OpenAI Research Catalog · Catalogue 052 · p. 7}}}
\expandafter\def\csname cite-universal\endcsname{\shortstack[l]{\refurl{https://github.com/openai/math/blob/1581d4cbdd69a0faf34b2089614fe2a814316a6a/preprints/Universal-cover-splitting-for-compact-Kahler-manifolds-September-23-2026/paper.pdf}{[OpenAI26a, Theorem 1.1, p. 3]}}}
\expandafter\def\csname cite-continuation\endcsname{\shortstack[l]{\refurl{https://github.com/openai/math/blob/1581d4cbdd69a0faf34b2089614fe2a814316a6a/preprints/Universal-cover-splitting-for-compact-Kahler-manifolds-September-23-2026/paper.pdf}{[OpenAI26a, Sections 5–6, pp. 12–17]}}}
\expandafter\def\csname cite-integrability\endcsname{\shortstack[l]{\refurl{https://github.com/openai/math/blob/1581d4cbdd69a0faf34b2089614fe2a814316a6a/preprints/Integrability-of-split-tangent-bundles-on-rationally-connected-manifolds-September-23-2026/main.pdf}{[OpenAI26b, Theorem 1.1, p. 1]}}}
\expandafter\def\csname cite-product\endcsname{\shortstack[l]{\refurl{https://github.com/openai/math/blob/1581d4cbdd69a0faf34b2089614fe2a814316a6a/preprints/Integrability-of-split-tangent-bundles-on-rationally-connected-manifolds-September-23-2026/main.pdf}{[OpenAI26b, Corollary 1.2, p. 2]}}}
\expandafter\def\csname cite-horing-product\endcsname{\shortstack[l]{\refurl{https://arxiv.org/pdf/math/0505327v3\string##page=2}{[Hor07, Theorem 1.4, p. 2]}}}
\expandafter\def\csname cite-companion\endcsname{\shortstack[l]{\refurl{https://github.com/openai/math/blob/1581d4cbdd69a0faf34b2089614fe2a814316a6a/preprints/Universal-cover-splitting-for-compact-Kahler-manifolds-September-23-2026/paper.pdf}{[OpenAI26a, Introduction, p. 3]}}}
\expandafter\def\csname cite-quotient\endcsname{\shortstack[l]{\refurl{https://github.com/openai/math/blob/1581d4cbdd69a0faf34b2089614fe2a814316a6a/preprints/Integrability-of-split-tangent-bundles-on-rationally-connected-manifolds-September-23-2026/main.pdf}{[OpenAI26b, Corollary 5.1, pp. 12–13]}}}

\Japanesetrue
\begin{document}
\begin{tikzpicture}
\node[half] (a) {\heading{}{Kählerの経路}{The Kähler route}
$X$ \Tx{コンパクト・連結Kähler}{compact connected Kähler}\\
$T_X=E_1\oplus E_2$,\quad $\operatorname{rk}E_i>0$\\
\Tx{両因子の可積分性を仮定}{Both summands assumed integrable}};
\node[half,right=.7cm of a] (b) {\heading{}{有理連結の経路}{The rationally connected route}
$X$ \Tx{滑らか・連結・射影・有理連結}{smooth, connected, projective, RC}\\
$T_X=E_1\oplus E_2$,\quad $\operatorname{rk}E_i>0$\\
\Tx{可積分性をまだ仮定しない}{Integrability not yet assumed}};
\node[half,below=2.4cm of a] (c) {\heading{}{全経路に沿う接続}{Continuation along all paths}
\Tx{正則boxから経路長方形へ}{Holomorphic boxes to path rectangles}\\
$H_o:L_1\times L_2\rightleftarrows X:K_o$\\
\Tx{互いに逆な局所芽}{Mutually inverse local germs}};
\node[half,below=2.4cm of b] (d) {\heading{}{可積分性を得る}{Obtain integrability}
$\Lambda^2E_i\longrightarrow T_X/E_i=E_j$\\
\Tx{Frobenius障害が消滅}{The Frobenius obstruction vanishes}};
\draw[edge] ($(a.south)+(-3,0)$)--($(c.north)+(-3,0)$);
\node[branchreason,anchor=west] at ($(a.south)!0.5!(c.north)+(-2.6,0)$) {\Tx{全polydiskを保つ延長と格子}{Extension and grids retaining full polydisks}\\\Cite{continuation}};
\draw[edge] ($(b.south)+(-3,0)$)--($(d.north)+(-3,0)$);
\node[branchreason,anchor=west] at ($(b.south)!0.5!(d.north)+(-2.6,0)$) {\Tx{tangent box・束の移送・二点評価}{Tangent boxes, transport and two-point evaluation}\\\Cite{integrability}};
\node[half,result,below=2.3cm of c,text width=7.5cm] (e) {\heading{}{普通普遍被覆の積}{Ordinary universal-cover product}
$\widetilde X\simeq Y_1\times Y_2$\\
$d\Phi(\pi^*E_i)=\operatorname{pr}_i^*T_{Y_i}$};
\node[half,result,below=2.3cm of d,text width=7.5cm] (f) {\heading{}{$X$自身の積}{A product of $X$ itself}
$X\simeq X_1\times X_2$\\
$E_i=\operatorname{pr}_i^*T_{X_i}$};
\draw[edge] ($(c.south)+(-3,0)$)--($(e.north)+(-3,0)$);
\node[branchreason,anchor=west] at ($(c.south)!0.5!(e.north)+(-2.6,0)$) {\Tx{monodromy・基点を合わせた持ち上げ}{Monodromy and based covering lifts}\\\Cite{universal}};
\draw[edge] ($(d.south)+(-3,0)$)--($(f.north)+(-3,0)$);
\node[branchreason,anchor=west] at ($(d.south)!0.5!(f.north)+(-2.6,0)$) {\Tx{有理連結性と可積分な因子から積へ}{RC and an integrable factor give a product}\\\Cite{horing-product}\\\Cite{product}};

\end{tikzpicture}
\end{document}
